\subsection{Matrices y parámetros}

Se tiene en principio, la matriz \textbf{B} para luego armar las matrices \textbf{G} y \textbf{H}.

%% Matriz B

A continuación se muestran las matrices \textbf{G} y \textbf{H}.

\[
\mathbf{G} =
\left[
\begin{array}{cccccccccccc|cccccccccccc}
1&0&0&0&0&0&0&0&0&0&0&0&
1&0&0&1&1&0&0&0&1&1&1&1\\
0&1&0&0&0&0&0&0&0&0&0&0&
0&1&0&0&1&1&1&0&0&1&1&1\\
0&0&1&0&0&0&0&0&0&0&0&0&
0&0&1&1&0&1&0&1&0&1&1&1\\
0&0&0&1&0&0&0&0&0&0&0&0&
1&0&1&1&1&1&1&0&0&0&1&0\\
0&0&0&0&1&0&0&0&0&0&0&0&
1&1&0&1&1&1&0&1&0&0&0&1\\
0&0&0&0&0&1&0&0&0&0&0&0&
0&1&1&1&1&1&0&0&1&1&0&0\\
0&0&0&0&0&0&1&0&0&0&0&0&
0&1&0&1&0&0&1&1&1&1&0&1\\
0&0&0&0&0&0&0&1&0&0&0&0&
0&0&1&0&1&0&1&1&1&1&1&0\\
0&0&0&0&0&0&0&0&1&0&0&0&
1&0&0&0&0&1&1&1&1&0&1&1\\
0&0&0&0&0&0&0&0&0&1&0&0&
1&1&1&0&0&1&1&1&0&1&0&0\\
0&0&0&0&0&0&0&0&0&0&1&0&
1&1&1&1&0&0&0&1&1&0&1&0\\
0&0&0&0&0&0&0&0&0&0&0&1&
1&1&1&0&1&0&1&0&1&0&0&1
\end{array}
\right]
\]

\[
\mathbf{H} =
\left[
\begin{array}{cccccccccccc|cccccccccccc}
1&0&0&1&1&0&0&0&1&1&1&1&
1&0&0&0&0&0&0&0&0&0&0&0\\
0&1&0&0&1&1&1&0&0&1&1&1&
0&1&0&0&0&0&0&0&0&0&0&0\\
0&0&1&1&0&1&0&1&0&1&1&1&
0&0&1&0&0&0&0&0&0&0&0&0\\
1&0&1&1&1&1&1&0&0&0&1&0&
0&0&0&1&0&0&0&0&0&0&0&0\\
1&1&0&1&1&1&0&1&0&0&0&1&
0&0&0&0&1&0&0&0&0&0&0&0\\
0&1&1&1&1&1&0&0&1&1&0&0&
0&0&0&0&0&1&0&0&0&0&0&0\\
0&1&0&1&0&0&1&1&1&1&0&1&
0&0&0&0&0&0&1&0&0&0&0&0\\
0&0&1&0&1&0&1&1&1&1&1&0&
0&0&0&0&0&0&0&1&0&0&0&0\\
1&0&0&0&0&1&1&1&1&0&1&1&
0&0&0&0&0&0&0&0&1&0&0&0\\
1&1&1&0&0&1&1&1&0&1&0&0&
0&0&0&0&0&0&0&0&0&1&0&0\\
1&1&1&1&0&0&0&1&1&0&1&0&
0&0&0&0&0&0&0&0&0&0&1&0\\
1&1&1&0&1&0&1&0&1&0&0&1&
0&0&0&0&0&0&0&0&0&0&0&1
\end{array}
\right]
\]

En el archivo "reference.ods", se verifican las operaciones $G \cdot H^{T} = 0$ y $B^{2}=I$

\subsection{Dimensión y tasa}

Respecto a la dimensión tenemos que la matriz \textbf{G} pertenece al campo de Galois $GF(2)^{12x24}$. De esta forma tenemos $ k = 12 $ bits de información y $ n = 24 $ bits totales, incluyendo redundancia.
El código Golay en este caso resulta ($n, k, d_{min}$)=(24,12,8).

La tasa se calcula como el cociente entre la cantidad de bits de información y los bits totales.

$$ R= \frac{k}{n} = \frac{12}{24} $$

$$ \boxed{R= 0.5} $$

Es decir que se tiene la misma cantidad de bits de información como de redundancia.

\subsection{Overhead}

La cantidad de bits adiciones que se agregan para poder realizar la codificación en relación con la información original es la siguiente.
Se tienen los bits de información ($k$) y los bits en la palabra codificada ($n$). Para el código de Golay, en este caso se tiene:

$$ overhead = n - k = 24 \, bits - 12 \, bits $$

$$ \boxed{overhead = 12 \, bits} $$

\subsection{Codificación de información de prueba}

Se hace la codificación de la siguiente palabra:

$$ m = 0xA5C $$

La palabra codificada resulta en:

$$ \boxed{v = 0xA5C9A5} $$

Luego multiplicando la palabra codificada $v$ por $\textbf{H}^{T}$ se tiene:

$$ v \cdot \textbf{H}^{T} = 0 $$

Todo el procedimiento de cálculo manual se muestra en el archivo $reference.ods$.
